In kinetic reduced magnetohydrodynamics (KRMHD) the compressive fluctuations are a kinetic field carried by the Alfv\'enic turbulence. Their free energy moves to fine structure in parallel velocity by phase mixing and to small perpendicular scales by the turbulence, and we ask what sets their Hermite spectrum $W(m)$. The compressive sector is exactly a continuum of passive scalars, or slices, labelled by the parallel phase velocity. Each slice is carried by the $\boldsymbol E\times\boldsymbol B$ flow and streams along the perturbed magnetic field at its own speed, and when the induction equation is ideal it streams along field lines that are frozen into the flow. The nonlinearity therefore moves no free energy between phase velocities and none between field-aligned wavenumbers. We propose that its only effect on $W(m)$ is to remove free energy from every Hermite moment at the same rate $1/T$, the rate at which the turbulence carries each slice to the dissipation scale, while phase mixing proceeds at the field-aligned wavenumber $\kappa$ at which the free energy was injected. This gives $W(m)\propto m^{-1/2}\exp[-(2m)^{1/2}/(s\kappa T)]$, with $s=\beta_i^{1/2}$, and a break at $m_*=(s\kappa T)^2/2\sim\beta_i/\chi^2$, where $\chi$ is the nonlinearity parameter. In forced KRMHD turbulence at $\beta_i=1$ and $\chi\approx1$, the perpendicular dissipation of each Hermite moment is proportional to $W(m)$ to within 20\% for $m\le16$. With $T$ and the distribution of $\kappa$ measured independently, the model reproduces $W(m)$ to within 8\% for $m\le30$ and the field-aligned Hermite flux to within 7\% for $m\le24$, and neither falls below it at large $m$, as a return flux would require. A split of the flux by the sign of $k_z$ reports a backward flux of about 12\% of the forward flux, which the field-aligned flux does not have. \placeholder{one or two sentences on the $\beta_i$ scan at fixed $\chi$ ($\beta_i=0.1$, 0.3, 1): whether $m_*$ scales with $\beta_i$ and the break enters the hierarchy} \placeholder{one sentence on the dependence of $T$ on resolution}
